% Translated from: thesis/chapter_01_绪论.md
% Target journal: General
% Generated: 2026-04-03

% --- Paragraph 1 ---
% # 第一章 绪论

\chapter{Introduction}

% --- Paragraph 2 ---
% ## 1.1 研究背景与意义

\section{Research Background and Significance}

% --- Paragraph 3 ---
% 通用博弈游戏（General Game Playing, GGP）旨在使智能体在运行时接收形式化规则描述，并在缺乏特定领域先验的条件下完成博弈推理、动作选择与策略生成。与围绕单一棋类手工设计特征、评估函数和搜索规则的传统博弈程序不同，GGP 更强调规则层、状态机层和决策层的统一性，其研究价值不仅体现在游戏智能本身，也体现在对通用推理、通用搜索与跨任务学习能力的验证上 \cite{genesereth2005general,genesereth2026general}。因此，GGP 一直被视为研究让智能体在陌生规则环境下快速形成有效决策能力的重要试验场。

General Game Playing (GGP) aims to enable an agent to receive formal rule descriptions at runtime and perform game reasoning, action selection, and strategy generation without relying on domain-specific priors. Unlike traditional game programs that are designed around handcrafted features, evaluation functions, and search rules for a single board game, GGP emphasizes unification across the rule layer, state-machine layer, and decision layer. Its value lies not only in game intelligence itself, but also in validating general reasoning, general search, and cross-task learning capabilities \cite{genesereth2005general,genesereth2026general}. Therefore, GGP has long been regarded as an important testbed for studying how agents can rapidly form effective decision-making ability in unfamiliar rule environments.

% --- Paragraph 4 ---
% 在现有 GGP 研究中，蒙特卡洛树搜索（Monte Carlo Tree Search, MCTS）因其不依赖强领域知识、能够在给定状态机上直接前向模拟而被广泛采用 \cite{browne2012survey}。尤其在规则未知但可执行的场景中，MCTS 能够通过选择、扩展、模拟和回传逐步构建搜索树，为智能体提供一种相对稳健的决策基础。进一步地，启发式 rollout、历史统计和神经网络评估等增强手段，也不断被引入 MCTS，以改善其在复杂博弈环境中的搜索效率与决策质量 \cite{finnsson2008simulation,anthony2017thinking,silver2018general}。与此同时，针对 GGP 的特点：游戏在运行时给定、离线调参难以精确适配，Sironi 等进一步提出了面向单个游戏在线调节搜索控制参数的自适应 MCTS 思路，这也从侧面说明了统一规则环境下搜索增强问题的现实复杂性 \cite{sironi2020selfadaptive}。

In existing GGP research, Monte Carlo Tree Search (MCTS) is widely adopted because it does not depend on strong domain knowledge and can perform direct forward simulation on a given state machine \cite{browne2012survey}. Especially in settings where rules are unknown in advance but executable, MCTS incrementally builds a search tree through selection, expansion, simulation, and backpropagation, providing a relatively robust decision foundation for agents. Furthermore, enhancement techniques such as heuristic rollouts, historical statistics, and neural-network-based evaluation have continuously been introduced into MCTS to improve search efficiency and decision quality in complex game environments \cite{finnsson2008simulation,anthony2017thinking,silver2018general}. At the same time, considering a key property of GGP, namely that games are specified at runtime and are difficult to tune precisely offline, Sironi et al. proposed an adaptive MCTS approach that adjusts search-control parameters online for a specific game. This also indicates the practical complexity of search enhancement under unified rule environments \cite{sironi2020selfadaptive}.

% --- Paragraph 5 ---
% 然而，GGP 中的搜索增强问题并不像固定棋类那样容易处理。首先，不同游戏在规则结构、行动空间和状态表示上的差异较大，使得面向单一游戏有效的启发式知识很难直接迁移。其次，随着研究对象从完全信息棋盘游戏扩展到隐藏信息、知识约束甚至通信行为等更复杂的规则场景，统一表示与统一推理的难度还会进一步上升 \cite{partridge2022hidden,zammit2023communication}。再次，纯搜索方法在有限决策时间下往往受到模拟开销限制，尤其当规则解释和状态转移本身已经消耗较多计算资源时，搜索预算会进一步被压缩。最后，神经价值方法虽然为搜索增强提供了新的可能，但在 GGP 场景下仍面临统一表示困难、监督构造粗糙、跨游戏泛化不足以及与搜索集成稳定性不强等问题 \cite{goldwaser2020deep,mcewan2022knowledge,maras2024fast}。这说明，在统一规则环境下，如何构建一个既能适配多种游戏、又能真实提升有限预算搜索表现的价值评估模块，仍然是一个值得研究的重要问题。

However, search enhancement in GGP is not as straightforward as in fixed game domains. First, substantial differences in rule structure, action space, and state representation across games make it difficult to transfer heuristics that work well for a single game. Second, as research moves from perfect-information board games to more complex rule settings involving hidden information, knowledge constraints, and even communication behavior, the difficulty of unified representation and unified reasoning further increases \cite{partridge2022hidden,zammit2023communication}. Third, pure search methods are often constrained by simulation overhead under limited decision time, and the search budget is further compressed when rule interpretation and state transition themselves already consume considerable computation. Finally, although neural value methods create new opportunities for search enhancement, they still face difficulties in GGP, including unified representation, coarse supervision design, insufficient cross-game generalization, and unstable integration with search \cite{goldwaser2020deep,mcewan2022knowledge,maras2024fast}. These points show that, in a unified rule environment, building a value-evaluation module that can both adapt to multiple games and genuinely improve search performance under limited budgets remains an important research problem.

% --- Paragraph 6 ---
% 从工程与应用角度看，这一问题具有双重意义。一方面，若能够在统一 GGP 框架中实现可复用的价值修正机制，就有望减少对特定游戏人工知识的依赖，提升通用智能体在新规则环境中的适应能力；另一方面，若价值模型能够真正转化为有限时间预算下的搜索收益，则说明“规则驱动执行 + 搜索蒸馏监督 + 残差价值修正”这一技术路线具备现实可行性。近期关于搜索、规划、学习与可解释性的专题讨论也持续将 GGP 相关问题列为重要方向，进一步说明该领域仍具有明显研究热度与拓展空间 \cite{soemers2024gametable}。因此，围绕 GGP 与 MCTS 开展搜索增强研究，不仅具有理论意义，也具有明确的系统实现与实验验证价值。

From an engineering and application perspective, this problem has dual significance. On one hand, if a reusable value-correction mechanism can be realized within a unified GGP framework, reliance on handcrafted knowledge for specific games may be reduced, thereby improving adaptation to new rule environments. On the other hand, if a value model can be translated into actual search gains under limited time budgets, it supports the practical feasibility of the technical route of rule-driven execution, search-distillation supervision, and residual value correction. Recent discussions on search, planning, learning, and interpretability have also continued to list GGP-related topics as important directions, indicating sustained research interest and room for extension in this area \cite{soemers2024gametable}. Therefore, search enhancement research around GGP and MCTS is important not only theoretically, but also for system implementation and experimental validation.

% --- Paragraph 7 ---
% ## 1.2 研究问题与挑战

\section{Research Questions and Challenges}

% --- Paragraph 8 ---
% 基于上述背景，本文聚焦的问题可概括为：在统一 GGP 框架下，如何利用搜索蒸馏与残差学习构建一个能够稳定接入 MCTS 的神经价值模型，使其在不破坏规则驱动执行逻辑的前提下，提高有限搜索预算下的叶节点评估质量和真实对局表现。

Based on the background above, this thesis focuses on the following question: under a unified GGP framework, how can search distillation and residual learning be used to build a neural value model that can be stably integrated into MCTS, so that the quality of leaf-node evaluation and real match performance can be improved under limited search budgets without breaking the logic of rule-driven execution.

% --- Paragraph 9 ---
% 这一问题至少包含三个层面的挑战。第一，GGP 的核心前提是规则在运行时给定，因此价值模型的输入表示必须尽量保持通用，不能依赖针对某一棋类的专用特征工程，否则模型很难跨游戏复用。第二，若仅以终局胜负或弱搜索回报作为监督，价值学习往往会面临监督噪声大、与真实搜索需求不匹配等问题，导致模型难以稳定学习到对搜索真正有帮助的评估信号。第三，即使训练出了价值模型，也仍然存在如何将其稳定接入 MCTS 的问题，即神经评估究竟应当替代什么、补充什么，以及如何在统一搜索骨架下验证其真实贡献。

This question involves at least three levels of challenges. First, the core premise of GGP is that rules are given at runtime, so the input representation of the value model must remain as general as possible and cannot rely on game-specific feature engineering; otherwise, cross-game reuse is difficult. Second, if only terminal win/loss outcomes or weak-search returns are used as supervision, value learning often suffers from high supervisory noise and mismatch with actual search demands, making it difficult for the model to learn evaluation signals that genuinely help search. Third, even after a value model is trained, stable integration into MCTS remains nontrivial: what should neural evaluation replace, what should it complement, and how can its real contribution be validated under a unified search skeleton.

% --- Paragraph 10 ---
% 结合已有章节内容可以看出，本文并不试图直接训练一个脱离搜索独立运行的通用博弈策略网络，而是将问题限定在价值修正这一更贴近当前项目条件与研究目标的范围内。也就是说，本文关注的不是完全重建整套决策能力，而是在已有低成本搜索估计的基础上，利用更强搜索生成 teacher 信号，再由模型学习修正量，使最终预测更接近强搜索评价，并进一步转化为在线搜索中的实战收益。这样的研究问题既保留了 GGP 的通用性要求，也与后文的系统实现和实验设计保持一致。

As indicated by the existing chapter content, this work does not attempt to directly train a general game policy network that operates independently of search. Instead, it narrows the problem to value correction, which is more aligned with current project conditions and research goals. In other words, rather than fully rebuilding the entire decision stack, this thesis starts from low-cost search estimates, uses stronger search to generate teacher signals, and then learns correction terms so that final predictions are closer to strong-search evaluations and can be converted into online playing strength. This formulation preserves the generality requirements of GGP while remaining consistent with the system implementation and experimental design in later chapters.

% --- Paragraph 11 ---
% ## 1.3 研究思路与主要工作

\section{Research Approach and Main Contributions}

% --- Paragraph 12 ---
% 围绕上述问题，本文提出 Search-Distilled Residual Progress-guided Value（SDRPV）方法，并据此构建了一套从规则执行、样本生成、模型训练到在线评测的完整实验闭环。整体思路是：首先在统一状态机上运行自对弈与搜索过程，为中间状态同时构造低成本基线值和高预算 teacher 值；随后采用基于棋盘 token 的统一状态表示与 Transformer 价值网络学习基线估计的修正量；最后将训练得到的价值模型接入共享 MCTS 内核，在固定时间预算下验证其相对于纯搜索、启发式搜索和神经价值基线的实际效果。

To address the above problem, this thesis proposes the Search-Distilled Residual Progress-guided Value (SDRPV) method and builds a complete experimental loop from rule execution and sample generation to model training and online evaluation. The overall idea is as follows: first, run self-play and search over a unified state machine to construct both low-cost baseline values and high-budget teacher values for intermediate states; second, use a unified board-token state representation and a Transformer value network to learn corrections to baseline estimates; finally, integrate the trained value model into a shared MCTS kernel and evaluate its practical effect against pure search, heuristic search, and neural-value baselines under a fixed-time budget.

% --- Paragraph 13 ---
% 本文的主要研究内容和工作可以概括为以下四个方面。

The main research content and contributions of this thesis can be summarized in four aspects.

% --- Paragraph 14 ---
% 1. 构建了面向 GGP 的统一规则执行与状态机基础。本文以 GDL/KIF 规则描述为输入，通过规则解析与逻辑推理链路实现统一状态机接口，使 `breakthrough`、`connectFour` 和 `hex` 三种游戏能够在一致的环境语义下完成合法动作查询、状态转移和终局评分，为后续方法比较提供统一执行基础。

1. It establishes unified rule execution and state-machine foundations for GGP. Using GDL/KIF rule descriptions as input, this work implements a unified state-machine interface through rule parsing and logical reasoning pipelines, so that `breakthrough`, `connectFour`, and `hex` can all perform legal-action queries, state transitions, and terminal scoring under a consistent environment semantics. This provides a unified execution basis for subsequent method comparison.

% --- Paragraph 15 ---
% 2. 提出了围绕搜索蒸馏与残差修正的 SDRPV 方法。本文不直接令模型回归完整价值，而是为同一状态同时构造低成本 baseline value 与高预算 teacher value，并以二者差值作为核心学习对象，从而将价值学习问题转化为稳定修正已有估计偏差的问题。这一设计使监督信号更贴近真实搜索需求，也使模型目标更贴近有限预算下的叶节点评估改进。

2. It proposes the SDRPV method centered on search distillation and residual correction. Instead of directly regressing the complete value, this work constructs both a low-cost baseline value and a high-budget teacher value for the same state and uses their difference as the core learning target, turning value learning into the more stable task of correcting existing estimation bias. This design makes supervision signals closer to real search demands and aligns the model objective with leaf-evaluation improvement under limited budgets.

% --- Paragraph 16 ---
% 3. 设计并实现了统一的状态编码、价值评估与搜索集成方案。本文采用基于棋盘 token 的通用表示方式，并以 Transformer 结构输出标量价值估计，使神经评估器能够在保持跨游戏统一性的同时融入共享 MCTS 内核。相较于直接替代规则系统的做法，本文始终保留规则真值在终局判定中的主导地位，从而明确了规则执行与神经评估的工程边界。

3. It designs and implements a unified scheme for state encoding, value estimation, and search integration. This work adopts a board-token-based general representation and outputs scalar value estimates with a Transformer architecture, allowing the neural evaluator to be integrated into a shared MCTS kernel while preserving cross-game consistency. Compared with directly replacing the rule system, this approach consistently preserves the dominant role of rule-grounded truth in terminal judgment, thereby clarifying the engineering boundary between rule execution and neural evaluation.

% --- Paragraph 17 ---
% 4. 完成了系统化的实验验证。本文以 `connectFour` 自对弈样本构建训练数据，并在 `breakthrough`、`connectFour` 和 `hex` 上进行 fixed-time（固定时间）在线评测。实验结果表明，SDRPV 在离线指标上优于基线值 `b` 与 teacher-only 路线，在在线 baseline benchmark 中对 `pure_mct` 和 `heuristic_mcts` 的三游戏总体胜率均达到 `0.7133`，对 `baseline_token_transformer` 也取得 `0.5467` 的总体正优势；ablation benchmark 则进一步说明，teacher 监督与 residual 目标对最终性能提升具有共同作用。上述结果说明，本文提出的技术路线能够在统一 GGP 框架下转化为真实的搜索收益。

4. It completes systematic experimental validation. This work uses `connectFour` self-play samples to build training data and performs fixed-time online evaluation on `breakthrough`, `connectFour`, and `hex`. Results show that SDRPV outperforms baseline value `b` and the teacher-only route on offline metrics. In the online baseline benchmark, SDRPV achieves an overall win rate of `0.7133` against both `pure_mct` and `heuristic_mcts` across the three games, and also obtains an overall positive edge of `0.5467` against `baseline_token_transformer`. The ablation benchmark further indicates that teacher supervision and the residual objective jointly contribute to final performance gains. These results demonstrate that the proposed technical route can be converted into real search gains within a unified GGP framework.

% --- Paragraph 18 ---
% 总体而言，本文的工作重点并不在于提出一种完全脱离既有搜索框架的新型博弈智能体，而在于围绕统一规则环境中的价值修正问题，形成一条具有清晰监督来源、明确系统边界和可验证实验收益的方法路线。这也是本文相对于一般神经博弈研究更强调通用规则场景下搜索增强的原因所在。
% 

Overall, the focus of this thesis is not to propose a brand-new game-playing agent that is fully detached from existing search frameworks. Instead, it targets value correction under a unified rule environment and develops a methodological route with clear supervision sources, explicit system boundaries, and verifiable experimental gains. This is also why, compared with general neural game-playing studies, this thesis places stronger emphasis on search enhancement in general-rule settings.

